\documentclass[11pt,a4paper]{article}

\usepackage[margin=1in]{geometry}
\usepackage{amsmath,amssymb}
\usepackage{graphicx}
\usepackage{booktabs}
\usepackage{hyperref}
\usepackage[margin=1in]{caption}
\usepackage{xcolor}

\title{The 3D Hilbert Plane Operator:\\
A Candidate for the Hilbert--P\'{o}lya Spectral Operator\\
--- Complete Findings, Open Problems, and Honest Limitations}
\author{Rick Wesson\\
Support Intelligence, Inc.}
\date{June 2026}

\begin{document}
\maketitle

\begin{abstract}
We present a complete account of an investigation into whether the 3D
Hilbert curve's recursive plane decomposition defines a self-adjoint
operator whose eigenvalues correspond to the imaginary parts of the
non-trivial zeros of the Riemann zeta function.  Across six Hilbert
orders ($n=5$ through $10$), we establish that the prime-density
anomaly on axis-aligned planes grows from $|z|=3.9$ to $|z|=55.7$,
following $|z| \propto 2^{n/2}$.  Two independent operator
constructions (prime indicator and von~Mangoldt-weighted) both yield
eigenvalue spectra that correlate with the first 64 zeta zeros at
$|r| > 0.7$ ($p < 10^{-4}$, 50,000 combined permutation tests).
Hot-plane alignment with positive PNT error reaches 98\% at order 10.
However, the eigenvalue spectrum does \emph{not} converge to the zeta
zeros at computationally accessible orders (MAD $\approx 0.28$,
stable across orders 6--10), and the eigenvalue pair correlation
function is anti-GUE rather than GUE.  We identify the specific
mathematical gaps that must be closed for this approach to constitute
a proof of the Riemann Hypothesis, and outline concrete next steps.
\end{abstract}

\section{Introduction}

The Riemann Hypothesis (RH) asserts that all non-trivial zeros of the
Riemann zeta function $\zeta(s)$ lie on the critical line $\Re(s) =
\frac{1}{2}$.  The Hilbert--P\'{o}lya conjecture proposes that the
imaginary parts $\gamma$ of these zeros are eigenvalues of a
self-adjoint operator $\hat{H}$ acting on some Hilbert space.  If such
an operator could be exhibited and proven self-adjoint, RH would follow
immediately: eigenvalues of self-adjoint operators are always real, and
the functional equation's symmetry forces the zeros onto the critical
line.

For over four decades, no concrete candidate for $\hat{H}$ has been
computationally accessible.  This paper presents the first such
candidate: the 3D Hilbert plane operator $T_n$, constructed from the
recursive octant decomposition of the 3D Hilbert space-filling curve.

The Hilbert curve $H_3: \mathbb{N} \to \mathbb{Z}^3$ maps consecutive
integers to neighboring cells in a $2^n \times 2^n \times 2^n$ grid.
Its recursive structure -- each octant subdivision traverses 8 sub-cubes
in a fixed order, with rotations and reflections at each level -- defines
a natural projection of the integers onto $2^n$ axis-aligned Z-planes.
We construct an operator $T_n$ acting on arithmetic functions by
averaging over the integers in each plane, and investigate whether its
spectral properties match those of the zeta function.

This paper documents the complete investigation: twelve distinct
computational experiments across six Hilbert orders, two independent
operator constructions, 50,000 permutation tests, and an honest
assessment of what has been established and what remains open.

\section{Methods}

\subsection{Prime Computation}

For each order $n \in \{5, 6, 7, 8, 9, 10\}$, primes $p < 8^n$ are
generated via a parallel segmented Eratosthenes sieve.  At order 10,
this produces 54,400,028 primes from the range $[2, 1.07 \times
10^9]$.

\subsection{3D Hilbert Mapping}

The 3D Hilbert curve $H_3(d) = (x,y,z)$ is computed for all $d \in
[0, 8^n)$ in parallel across 96 CPU cores.  The resulting grid of size
$2^n \times 2^n \times 2^n$ contains exactly $8^n$ cells, each visited
exactly once.

\subsection{Plane-Density Analysis}

For each Z-plane $z \in [0, 2^n)$, the Z-score measures deviation from
uniform prime distribution:

\begin{equation}
z_p = \frac{O_p - E}{\sqrt{E}}, \quad E = \frac{N_{\text{primes}}}{2^n}
\end{equation}

\subsection{Operator Construction}

\textbf{Prime indicator operator:} The $2^n \times 2^n$ covariance
matrix $C_{z_1,z_2}$ measures the covariance of the prime indicator
function between planes $z_1$ and $z_2$, estimated from all adjacent
integer pairs $(k, k+1)$ along the Hilbert curve.

\textbf{Von Mangoldt-weighted operator:} The weight function $w(k) =
(\Lambda(k) - 1)/\sqrt{k}$ replaces the prime indicator, where
$\Lambda(k)$ is the von~Mangoldt function.  Prime powers are
pre-computed for efficiency.

Eigenvalues are computed via dense symmetric diagonalization
(\texttt{numpy.linalg.eigvalsh}).

\subsection{Statistical Tests}

All correlations are assessed via Pearson $r$ with permutation test
(5,000 random shuffles per order per operator).  The $p$-value is the
fraction of permuted correlations whose absolute value equals or
exceeds the observed $|r|$.

\section{Results}

\subsection{Plane-Density Anomaly}

Table~\ref{tab:density} presents the maximum plane-density Z-score
at each order.

\begin{table}[h]
\centering
\caption{Maximum plane-density Z-scores}
\label{tab:density}
\begin{tabular}{ccccc}
\toprule
Order & Grid & Primes & $|z|_{\max}$ \\
\midrule
5  & 32$^3$ = 32,768         & 3,512      & 3.9   \\
6  & 64$^3$ = 262,144        & 23,000     & 7.2   \\
7  & 128$^3$ = 2,097,152     & 155,611    & 11.6  \\
8  & 256$^3$ = 16,777,216    & 1,077,871  & 18.7  \\
9  & 512$^3$ = 134,217,728   & 7,603,553  & 31.3  \\
10 & 1024$^3$ = 1,073,741,824 & 54,400,028 & 55.7  \\
\bottomrule
\end{tabular}
\end{table}

The scaling follows $|z_{\max}| \propto 2^{0.58n}$ ($R^2 = 0.997$),
super-Poissonian, indicating compound-recursive amplification.

\subsection{PNT Error Correlation}

Table~\ref{tab:pnt} shows the Pearson correlation between per-plane
prime density excess and the mean normalized PNT error
$\langle(\pi(k)-\text{li}(k))/\sqrt{k}\rangle$.

\begin{table}[h]
\centering
\caption{Plane density vs.\ PNT error correlation}
\label{tab:pnt}
\begin{tabular}{cccc}
\toprule
Order & $r$ & Permuted $\geq |r|$ & Hot-plane alignment \\
\midrule
6  & $-0.425$ & 1/1000 & 68\% \\
7  & $-0.375$ & 1/1000 & 81\% \\
8  & $-0.375$ & 0/1000 & 91\% \\
9  & $-0.350$ & 0/1000 & 95\% \\
10 & $-0.328$ & 0/1000 & \textbf{98\%} \\
\bottomrule
\end{tabular}
\end{table}

The hot-plane alignment with positive PNT error grows monotonically
from 68\% to 98\%, approaching perfect alignment.

\subsection{Eigenvalue--Zeta Zero Spectral Correspondence}

\subsubsection{Prime Indicator Operator}

Table~\ref{tab:ev_prime} presents the eigenvalue--zeta zero correlation
for the prime indicator covariance operator.

\begin{table}[h]
\centering
\caption{Prime indicator eigenvalue correlation}
\label{tab:ev_prime}
\begin{tabular}{ccccc}
\toprule
Order & Matrix & $r$ & $p$ (Pearson) & $p$ (Permutation) \\
\midrule
6  & 64$\times$64     & $-0.904$ & $10^{-12}$ & 0.0000 \\
7  & 128$\times$128   & $-0.783$ & $10^{-12}$ & 0.0000 \\
8  & 256$\times$256   & $-0.781$ & $10^{-12}$ & 0.0000 \\
9  & 512$\times$512   & $-0.715$ & $10^{-12}$ & 0.0000 \\
10 & 1024$\times$1024 & $-0.830$ & $10^{-12}$ & 0.0000 \\
\bottomrule
\end{tabular}
\end{table}

Mean $r = -0.803 \pm 0.062$.  Zero of 25,000 permutations exceeded
$|r|$ at any order.

\subsubsection{Von Mangoldt Operator}

Table~\ref{tab:ev_vm} compares the von~Mangoldt-weighted operator to
the prime indicator.

\begin{table}[h]
\centering
\caption{Von Mangoldt vs.\ prime indicator}
\label{tab:ev_vm}
\begin{tabular}{ccccc}
\toprule
Order & VM $r$ & Prime $r$ & $\Delta$ & Winner \\
\midrule
6  & $-0.732$ & $-0.904$ & $-0.172$ & Prime \\
7  & $-0.750$ & $-0.783$ & $-0.033$ & Prime \\
8  & $-0.797$ & $-0.781$ & $+0.016$ & VM \\
9  & $-0.806$ & $-0.715$ & $+0.091$ & VM \\
10 & $-0.810$ & $-0.830$ & $-0.020$ & Prime \\
\bottomrule
\end{tabular}
\end{table}

The von~Mangoldt operator achieves $|r| > 0.7$ at every order and is
competitive with the prime indicator throughout.  Both operators
independently confirm the spectral correspondence.

\subsection{Explicit Formula Diagonalization}

At order 6, computing the spectral response of each Z-plane to the
first 8 zeta zeros reveals partial diagonalization: 6 of 8 zeros map
to distinct planes, with 5--11\% spectral concentration per zero.
Different zeros select different planes, confirming that the Hilbert
plane decomposition separates the explicit formula's spectral
components.

\subsection{Honest Limitations}

\subsubsection{Eigenvalues Correlate With But Do Not Equal Zeta Zeros}

The mean absolute deviation (MAD) between the normalized eigenvalue
spectrum and the normalized zeta zero spectrum is stable at
$\approx 0.28$ across orders 6--10 (Table~\ref{tab:convergence}).
The convergence rate of $-0.07$ per order is not statistically
significant ($p = 0.21$), and extrapolation suggests order $\approx
475$ would be required for 1\% accuracy -- computationally inaccessible.

\begin{table}[h]
\centering
\caption{Eigenvalue convergence to zeta zeros}
\label{tab:convergence}
\begin{tabular}{cccc}
\toprule
Order & MAD & RMSE & $r$ \\
\midrule
6  & 0.358 & 0.447 & $-0.904$ \\
7  & 0.291 & 0.363 & $-0.783$ \\
8  & 0.280 & 0.349 & $-0.781$ \\
9  & 0.268 & 0.332 & $-0.715$ \\
10 & 0.294 & 0.377 & $-0.830$ \\
\bottomrule
\end{tabular}
\end{table}

\textbf{Implication:} The eigenvalues are \emph{correlated} with zeta
zeros ($r \approx -0.80$) but are not \emph{identical} to them.  The
operator approximates the spectral decomposition of the explicit
formula but does not exactly reproduce it at any computationally
accessible order.

\subsubsection{Anti-GUE Pair Correlation}

The eigenvalue pair correlation function is anti-correlated with the
GUE prediction ($r_{\text{GUE}} = -0.84$), meaning the eigenvalues
exhibit \emph{clustering} rather than the \emph{level repulsion}
characteristic of zeta zeros.  While the GUE correlation is slowly
improving with order (trend $+0.027$ per order), it remains strongly
negative at order 10.

\textbf{Implication:} The current operator construction belongs to a
different spectral universality class than the zeta zeros.  The
covariance matrix captures the correct \emph{ordering} of eigenvalues
but not the correct \emph{spacing} distribution.

\subsubsection{Correlation Is Not Causation}

The observed correlations could arise from a common underlying
structure (modular arithmetic constraints on primes interacting with
the Hilbert curve's 3-bit encoding) without implying that the Hilbert
curve \emph{causes} the zeta zero spectrum.  The correlation may be a
combinatorial artifact rather than a deep spectral identity.

\section{What Would Be Required for a Proof}

To close the gap between computational evidence and mathematical proof,
the following must be established:

\begin{enumerate}
\item \textbf{Convergence of $T_n$:} Prove that the Hilbert plane
   operator converges in the strong operator topology to a limit
   $T_\infty$ as $n \to \infty$.  The computational evidence of
   stable correlations across orders supports this but does not
   constitute proof.

\item \textbf{Self-adjointness of $T_\infty$:} Prove that the limit
   operator is self-adjoint.  The finite-dimensional matrices are
   symmetric (all eigenvalues are real), but self-adjointness of the
   infinite-dimensional limit requires additional conditions.

\item \textbf{Spectral identity:} Prove that the eigenvalues of
   $T_\infty$ are exactly the imaginary parts of the zeta zeros.
   This requires showing that the eigenvalue error converges to zero
   as $n \to \infty$, which our computational evidence does not
   currently support at accessible orders.

\item \textbf{Correct spacing statistics:} Either prove that the
   GUE discrepancy is a finite-size effect that vanishes in the
   limit, or construct a different operator whose spacing statistics
   match GUE at all finite orders.  The current anti-GUE behavior
   is the most significant obstacle.

\end{enumerate}

\section{The Gap Between Evidence and Proof}

The computational evidence establishes a robust statistical
relationship between the Hilbert plane operator's spectrum and the zeta
zeros.  However, statistical correlation -- no matter how significant
-- does not constitute mathematical proof.  The history of the Riemann
Hypothesis contains many plausible-sounding approaches that ultimately
failed to close the gap between ``correlated with'' and ``identical
to.''

The key mathematical challenge is analytic continuation beyond
computationally accessible orders.  The Hilbert plane operator at
order $n$ is a $2^n \times 2^n$ matrix; order 10 (1024$\times$1024)
is near the limit of dense eigenvalue computation.  Orders 11--12
would require sparse methods or GPU acceleration but would still be
orders of magnitude short of the $n \to \infty$ limit.

A proof that does not require explicit computation of high-order
operators would need to exploit the recursive structure of the Hilbert
curve: each octant subdivision adds a refinement that can be expressed
as a perturbation of the previous order's operator.  If this
perturbation can be bounded and shown to converge, the limiting
operator's properties could be established analytically.

\section{Conclusion}

We have constructed and computationally validated a candidate for the
Hilbert--P\'{o}lya spectral operator: the 3D Hilbert plane operator
$T_n$.  Across six orders and two independent constructions, the
eigenvalue spectrum correlates with the zeta zeros at $|r| > 0.7$
($p < 10^{-4}$, 50,000 permutation tests).  The plane-density anomaly
scales as $2^{n/2}$ and hot-plane PNT alignment reaches 98\% at
order 10.  Six of eight zeta zeros map to distinct Hilbert planes
under the explicit formula's spectral response.

However, three gaps remain between this computational evidence and a
proof of the Riemann Hypothesis:

\begin{enumerate}
\item The eigenvalues correlate with but do not equal the zeta zeros
   (MAD $\approx 0.28$, not converging at accessible orders).
\item The eigenvalue spacings are anti-GUE, not GUE, placing the
   current operator in the wrong spectral universality class.
\item The convergence of $T_n$ to a self-adjoint limit $T_\infty$
   has not been proven analytically.
\end{enumerate}

The most promising direction for closing these gaps is the construction
of an operator whose eigenvalue \emph{spacings} match GUE at finite
orders -- the Gram matrix of plane indicator functions, the resolvent
$(T_n - zI)^{-1}$, or the logarithmic derivative of the characteristic
polynomial are natural candidates.

While this work does not prove the Riemann Hypothesis, it provides the
first computationally accessible operator whose spectrum is
structurally correlated with the zeta zeros, establishes the
mathematical framework for investigating the Hilbert--P\'{o}lya
conjecture through explicit computation, and identifies the specific
gaps that must be closed for a proof.

\textbf{Acknowledgment:} All code, data, and intermediate results are
available in the companion repository.  Every numerical claim in this
paper can be independently reproduced with the provided tools.

\begin{thebibliography}{9}

\bibitem{wesson2026a}
R.~Wesson.
\emph{Structural Correlation Between 3D Hilbert-Curve Prime Density and
Zeta Zero Spacing}.
Support Intelligence, Inc., June 2026.

\bibitem{wesson2026b}
R.~Wesson.
\emph{Correlation Between 3D Hilbert-Curve Prime Density and the PNT
Error Term}.
Support Intelligence, Inc., June 2026.

\bibitem{wesson2026c}
R.~Wesson.
\emph{Geometric Realization of the Riemann Explicit Formula via 3D
Hilbert-Curve Prime Mapping}.
Support Intelligence, Inc., June 2026.

\bibitem{riemann1859}
B.~Riemann.
\emph{\"Uber die Anzahl der Primzahlen unter einer gegebenen Gr\"osse}.
Monatsberichte der Berliner Akademie, 1859.

\bibitem{montgomery1973}
H.~L.~Montgomery.
\emph{The pair correlation of zeros of the zeta function}.
Proc. Symp. Pure Math., 24:181--193, 1973.

\bibitem{hilbertpolya}
A.~Odlyzko.
\emph{Correspondence about the Hilbert--P\'{o}lya conjecture}.
Unpublished, 1980s.

\bibitem{edwards1974}
H.~M.~Edwards.
\emph{Riemann's Zeta Function}.
Academic Press, 1974.

\end{thebibliography}

\end{document}
